\documentclass[12pt]{article}

% =========================
% Toggle: show only problems, or problems + student solution boxes
% =========================
\newif\ifsolutions
\solutionstrue     % Use this for the student template with solution boxes.
%\solutionsfalse   % Use this for a question sheet only.

% =========================
% Packages
% =========================
\usepackage[a4paper,margin=24mm]{geometry}
\usepackage{amsmath,amssymb,mathtools}
\usepackage{braket}
\usepackage{siunitx}
\usepackage[most]{tcolorbox}
\usepackage{enumitem}
\usepackage{graphicx}
\usepackage{xcolor}
\usepackage[T1]{fontenc}
\usepackage{lmodern}
\usepackage{microtype}
\usepackage{float} % Allows figure to be placed at exact desired location
\usepackage[hidelinks]{hyperref}

\sisetup{inter-unit-product = \ensuremath{{}\cdot{}}}
\renewcommand{\d}{\mathop{}\!\mathrm{d}}
\setlength{\parindent}{0pt}
\setlength{\parskip}{0.55em}
\setlist[enumerate,1]{label=(\alph*),itemsep=0.25em,topsep=0.25em}
\setlist[enumerate,2]{label=(\roman*),itemsep=0.15em,topsep=0.15em}

% =========================
% Boxes
% =========================
\tcbset{
  enhanced,
  breakable,
  boxrule=0.8pt,
  arc=1.5mm,
  left=2mm,right=2mm,top=1.5mm,bottom=1.5mm
}

\newtcolorbox{problembox}[1]{
  colback=gray!8,
  colframe=black,
  fonttitle=\bfseries,
  title={#1}
}

\newtcolorbox{solutionbox}{
  colback=blue!3,
  colframe=blue!60!black,
  title={Solution},
  fonttitle=\bfseries
}

\newtcolorbox{answerbox}{
  colback=yellow!10,
  colframe=orange!70!black,
  title={Final answer},
  fonttitle=\bfseries
}

\newtcolorbox{insightbox}{
  colback=green!5,
  colframe=green!50!black,
  title={Physical insight},
  fonttitle=\bfseries
}

\newcommand{\problemsep}{\vspace{0.8em}}

\newcommand{\placeholderfigure}[2][50mm]{%
\begin{center}
\fbox{\parbox[c][#1][c]{0.82\linewidth}{\centering\small Figure placeholder}}\\[0.3em]
{\footnotesize\emph{Caption: #2}}
\end{center}}

\graphicspath{{figures/}} % If you want to place figures in the folder 'figures' (cleaner)

\setcounter{secnumdepth}{0} % Avoids numbering of the section and subsection while allowing them to appear in the TOC



% ======================================================
% ======================================================
% ================                    ==================
% ================   BEGIN DOCUMENT   ==================
% ================                    ==================
% ======================================================
% ======================================================

\begin{document}

\begin{center}
{\Large\bfseries QT01 -- Foundations of Quantum Mechanics\\[0.2em]
Assessed problem set (PRB) -- Week 4\\[0.3em]
{\normalsize Weight: 30\% of module grade}\\[1em]
Candidate number: 123456  % Replace candidate number with yours
}
\end{center}


% ======================================================
%                  Question 1
% ======================================================
\section{Question 1}
\begin{problembox}{Question 1 \hfill [8 marks]}
Consider the operator
\[
    \hat{A} = \begin{pmatrix} -3 & 2i \\ -2i & -3 \end{pmatrix}.
\]
Prove that \( \hat{A} \) is Hermitian.
\end{problembox}

% ======================================================
%                  Solution 1
% ======================================================
\ifsolutions
\begin{solutionbox}
 Replace this placeholder with your own detailed derivation, making your algebra clear and showing the steps needed to reach the result.
\end{solutionbox}

\begin{answerbox}
State your final answer(s) clearly here.
\end{answerbox}

\begin{insightbox}
Add a short physical interpretation here. Explain what the result means, why it matters, or what feature of the system it illustrates.
\end{insightbox}
\fi

\newpage
% ======================================================
%                  Question 2
% ======================================================
\section{Question 2}
\begin{problembox}{Question 2 \hfill [37 marks]}
The wave function of a particle moving in one dimension is given by
\[
    \psi(x) = \begin{cases}
        B \exp(\beta x) & \text{for } x<0,\\
        B \exp(-2\beta x) & \text{for } x \ge 0,
    \end{cases}
\]
where \(\beta\) is a real and positive constant.

\begin{enumerate}
    \item What is the probability density for finding the particle along the \(x\)-axis? Sketch the probability density function. \hfill [7 marks]
    \item Calculate the normalisation constant \(B\). \hfill [10 marks]
    \item Calculate the average position of the particle on the \(x\)-axis, as a function of \(\beta\). Mark the average position in your sketch. \hfill [15 marks]
    \item Mark the most likely position of the particle in your sketch. \hfill [5 marks]
\end{enumerate}

\textit{Hint: You can calculate the integrals you need by expressing powers of \(x\) through repeated differentiation with respect to the parameter in the exponential, for example}
\begin{align*}
    \int_a^b x \exp(\gamma x)\,\d x
    &= \frac{\partial}{\partial \gamma} \int_a^b \exp(\gamma x)\,\d x, \\
    \int_a^b x^2 \exp(\gamma x)\,\d x
    &= \frac{\partial^2}{\partial \gamma^2} \int_a^b \exp(\gamma x)\,\d x,
\end{align*}
\textit{and so on.}
\end{problembox}

% ======================================================
%                  Solution 2
% ======================================================
\ifsolutions
\begin{solutionbox}
 Replace this placeholder with your own detailed derivation, making your algebra clear and showing the steps needed to reach the result.
\end{solutionbox}

\begin{answerbox}
State your final answer(s) clearly here.
\end{answerbox}

\begin{insightbox}
Add a short physical interpretation here. Explain what the result means, why it matters, or what feature of the system it illustrates.
\end{insightbox}
\fi

\newpage
% ======================================================
%                  Question 3
% ======================================================
\section{Question 3}
\begin{problembox}{Question 3 \hfill [22 marks]}
Take the wave function
\[
    \varphi(x) = N \sin(kx).
\]
\begin{enumerate}
    \item Calculate the effect of the operator
    \[
        \left[x\left(\frac{\hbar}{i}\frac{\partial}{\partial x}\right)\right]
    \]
    on \(\varphi(x)\). \hfill [2 marks]

    \item Calculate the effect of the operator
    \[
        \left[\left(\frac{\hbar}{i}\frac{\partial}{\partial x}\right)x\right]
    \]
    on \(\varphi(x)\), that is, multiply \(\varphi(x)\) by \(x\) before acting with the differential operator.
    \\
    \phantom{add some invisible space} \hfill [5 marks]

    \item Calculate the difference of the two previous results, that is,
    \[
        \text{result of (a)} - \text{result of (b)}.
    \]
    Express the answer in terms of \(\varphi(x)\). \hfill [5 marks]

    \item Repeat parts (a) to (c) for a general wave function \(\psi(x)\), that is, calculate
    \[
        \left[x\left(\frac{\hbar}{i}\frac{\partial}{\partial x}\right)\right] \psi(x)
        -
        \left[\left(\frac{\hbar}{i}\frac{\partial}{\partial x}\right)x\right] \psi(x).
    \]
    \hfill [10 marks]
\end{enumerate}
\end{problembox}

% ======================================================
%                  Solution 3
% ======================================================
\ifsolutions
\begin{solutionbox}
 Replace this placeholder with your own detailed derivation, making your algebra clear and showing the steps needed to reach the result.
\end{solutionbox}

\begin{answerbox}
State your final answer(s) clearly here.
\end{answerbox}

\begin{insightbox}
Add a short physical interpretation here. Explain what the result means, why it matters, or what feature of the system it illustrates.
\end{insightbox}
\fi

\newpage
% ======================================================
%                  Question 4
% ======================================================
\section{Question 4}
\begin{problembox}{Question 4 \hfill [8 marks]}
A system is prepared in a state
\[
    \ket{\psi} = \frac{1}{\sqrt{3}} \ket{\phi_1}
    - \frac{\sqrt{2}}{\sqrt{3}} \ket{\phi_2}
    + \frac{i}{\sqrt{3}} \ket{\phi_3},
\]
where \( \ket{\phi_1} \), \( \ket{\phi_2} \), and \( \ket{\phi_3} \) are orthonormal eigenstates of an observable \(A\), with corresponding eigenvalues \(a_1\), \(a_2\), and \(a_3\), respectively. What is the probability of measuring \(A\) to be \(a_3\)?
\end{problembox}

% ======================================================
%                  Solution 4
% ======================================================
\ifsolutions
\begin{solutionbox}
 Replace this placeholder with your own detailed derivation, making your algebra clear and showing the steps needed to reach the result.
\end{solutionbox}

\begin{answerbox}
State your final answer(s) clearly here.
\end{answerbox}

\begin{insightbox}
Add a short physical interpretation here. Explain what the result means, why it matters, or what feature of the system it illustrates.
\end{insightbox}
\fi

\newpage
% ======================================================
%                  Question 5
% ======================================================
\section{Question 5}
\begin{problembox}{Question 5 \hfill [25 marks]}
A system is initially in the state
\[
    \ket{\psi_0} = \frac{1}{\sqrt{7}}\left[\sqrt{2}\ket{\phi_1}
    + \sqrt{3}\ket{\phi_2}
    + \ket{\phi_3}
    + \ket{\phi_4} \right],
\]
where \(\ket{\phi_n}\) are eigenstates of the system's Hamiltonian such that
\[
    \hat{H}\ket{\phi_n} = n^2\mathcal{E}_0\ket{\phi_n}.
\]

\begin{enumerate}
    \item If energy is measured, what values will be obtained and with what probabilities? \\
     \phantom{add some invisible space}\hfill [10 marks]

    \item Consider an operator \(\hat{A}\) whose action on \(\ket{\phi_n}\) is defined by
    \[
        \hat{A}\ket{\phi_n} = (n+1)a_0\ket{\phi_n}.
    \]
    If \(A\) is measured, what values will be obtained and with what probabilities? \\
    \phantom{add some invisible space}\hfill [10 marks]

    \item Suppose that a measurement of the energy yields \(4\mathcal{E}_0\). If we measure \(A\) immediately afterwards, what value will be obtained? \hfill [5 marks]
\end{enumerate}
\end{problembox}

% ======================================================
%                  Solution 5
% ======================================================
\ifsolutions
\begin{solutionbox}
 Replace this placeholder with your own detailed derivation, making your algebra clear and showing the steps needed to reach the result.
\end{solutionbox}

\begin{answerbox}
State your final answer(s) clearly here.
\end{answerbox}

\begin{insightbox}
Add a short physical interpretation here. Explain what the result means, why it matters, or what feature of the system it illustrates.
\end{insightbox}
\fi

\end{document}
